% Part: incompleteness
% Chapter: theories-computability
% Section: complete-decidable

\documentclass[../../../include/open-logic-section]{subfiles}

\begin{document}

\olfileid{inc}{tcp}{cdc}

\olsection{\printtoken{S}{axiomatizable} પૂર્ણ સિદ્ધાંતો નિર્ણેય છે}

સિદ્ધાંતને {\em પૂર્ણ} ત્યારે કહીએ
જ્યારે દરેક વાક્ય~$!A$ માટે
$!A$ અથવા $\lnot !A$
બેમાંથી એક સાબિત થાય.

\begin{lem}
માનો કે સિદ્ધાંત~$\Th{T}$
પૂર્ણ અને !!{axiomatizable} છે.
તો $\Th{T}$ નિર્ણેય છે.
\end{lem}

\begin{proof}
માનો કે $\Th{T}$ પૂર્ણ છે
અને $A$ તેના સ્વયંસિદ્ધિઓનો
સંગણનીય ગણ છે.
જો $\Th{T}$ અસુસંગત હોય,
તો તે સ્પષ્ટપણે સંગણનીય છે.
(કલનવિધિ: ``હા'' કહો.)
આથી માની શકીએ કે
$\Th{T}$ સુસંગત પણ છે.

વાક્ય~$!A$ એ~$\Th{T}$માં
છે કે નહીં તે નક્કી કરવા,
$\Th{T}$માંથી~$!A$નું
!!a{derivation} અને
$\lnot !A$નું !!a{derivation}
બંને એકસાથે શોધો.
$\Th{T}$ પૂર્ણ હોવાથી
બેમાંથી એક મળી જ જશે;
અને $\Th{T}$ સુસંગત હોવાથી,
જો $\lnot !A$નું
!!a{derivation} મળે તો
$!A$નું !!{derivation}
હોઈ શકે નહીં.

જુદી રીતે કહીએ તો
$\Th{T}$ !!{c.e.} છે
તે આપણે જાણીએ છીએ;
અગાઉ સાબિત કરેલા પ્રમેય
મુજબ $\Th{T}$નો પૂરક પણ
!!{c.e.} છે એટલું
બતાવવું પૂરતું છે.
પરંતુ !!a{formula}~$!A$ એ~$\Th{\bar
T}$માં ત્યારે અને માત્ર
ત્યારે જ હોય જ્યારે
$\lnot !A$ એ~$\Th{T}$માં
હોય; તેથી
$\Th{\bar T} \leq_m
\Th{T}$.
\end{proof}

\end{document}
